\documentclass[12pt, a4paper]{article}

% ─────────────────────────────────────────────────────────────
% Packages
% ─────────────────────────────────────────────────────────────
\usepackage[utf8]{inputenc}
\usepackage[T1]{fontenc}
\usepackage[french]{babel}
\usepackage{lmodern}
\usepackage{geometry}
\usepackage{graphicx}
\usepackage{xcolor}
\usepackage{hyperref}
\usepackage{booktabs}
\usepackage{tabularx}
\usepackage{multirow}
\usepackage{array}
\usepackage{colortbl}
\usepackage{listings}
\usepackage{fancyhdr}
\usepackage{titlesec}
\usepackage{titling}
\usepackage{enumitem}
\usepackage{mdframed}
\usepackage{tcolorbox}
\usepackage{amsmath}
\usepackage{amssymb}
\usepackage{float}
\usepackage{caption}
\usepackage{subcaption}
\usepackage{setspace}
\usepackage{microtype}
\usepackage{parskip}
\usepackage{tikz}
\usepackage{pgfplots}
\pgfplotsset{compat=1.18}
\usetikzlibrary{shapes.geometric, arrows.meta, positioning, fit, backgrounds}

% ─────────────────────────────────────────────────────────────
% Mise en page
% ─────────────────────────────────────────────────────────────
\geometry{
  top=2.5cm, bottom=2.5cm,
  left=2.5cm, right=2.5cm,
  headheight=15pt
}

% ─────────────────────────────────────────────────────────────
% Couleurs
% ─────────────────────────────────────────────────────────────
\definecolor{navy}{HTML}{0D1B2A}
\definecolor{blue}{HTML}{1565C0}
\definecolor{lblue}{HTML}{1E88E5}
\definecolor{accent}{HTML}{00ACC1}
\definecolor{silver}{HTML}{ECEFF1}
\definecolor{dgray}{HTML}{37474F}
\definecolor{mgray}{HTML}{78909C}
\definecolor{green}{HTML}{2E7D32}
\definecolor{orange}{HTML}{E65100}
\definecolor{codebg}{HTML}{F3F4F6}
\definecolor{codeframe}{HTML}{D1D5DB}

% ─────────────────────────────────────────────────────────────
% Listings (code)
% ─────────────────────────────────────────────────────────────
\lstset{
  backgroundcolor=\color{codebg},
  basicstyle=\ttfamily\small,
  keywordstyle=\color{blue}\bfseries,
  commentstyle=\color{mgray}\itshape,
  stringstyle=\color{green},
  frame=single,
  rulecolor=\color{codeframe},
  framexleftmargin=8pt,
  xleftmargin=12pt,
  numbers=left,
  numberstyle=\tiny\color{mgray},
  numbersep=8pt,
  breaklines=true,
  breakatwhitespace=true,
  tabsize=4,
  showstringspaces=false,
  captionpos=b,
  language=Python,
  literate={é}{{\'e}}1 {è}{{\`e}}1 {ê}{{\^e}}1 {à}{{\`a}}1 {ù}{{\`u}}1
           {ç}{{\c c}}1 {î}{{\^i}}1 {ô}{{\^o}}1 {û}{{\^u}}1
           {É}{{\'E}}1 {È}{{\`E}}1 {À}{{\`A}}1
           {→}{{$\rightarrow$}}1
           {≈}{{$\approx$}}1 {≤}{{$\leq$}}1 {≥}{{$\geq$}}1
}

% ─────────────────────────────────────────────────────────────
% tcolorbox styles
% ─────────────────────────────────────────────────────────────
\tcbuselibrary{most}

\newtcolorbox{infobox}[1][]{
  colback=blue!5!white, colframe=blue!60!white,
  fonttitle=\bfseries\small, title=#1,
  arc=4pt, boxrule=0.6pt,
  left=8pt, right=8pt, top=6pt, bottom=6pt
}

\newtcolorbox{resultbox}{
  colback=green!5!white, colframe=green!50!white,
  arc=4pt, boxrule=0.6pt,
  left=8pt, right=8pt, top=6pt, bottom=6pt
}

\newtcolorbox{warningbox}{
  colback=orange!8!white, colframe=orange!60!white,
  arc=4pt, boxrule=0.6pt,
  left=8pt, right=8pt, top=6pt, bottom=6pt
}

% ─────────────────────────────────────────────────────────────
% En-têtes et pieds de page
% ─────────────────────────────────────────────────────────────
\pagestyle{fancy}
\fancyhf{}
\renewcommand{\headrulewidth}{0.4pt}
\renewcommand{\footrulewidth}{0.4pt}
\fancyhead[L]{\small\textcolor{mgray}{S2R Project --- INPT}}
\fancyhead[R]{\small\textcolor{mgray}{Rapport Technique Complet}}
\fancyfoot[C]{\small\textcolor{mgray}{\thepage}}
\fancyfoot[L]{\small\textcolor{mgray}{Institut National des Postes et Telecommunications}}

% ─────────────────────────────────────────────────────────────
% Titres des sections
% ─────────────────────────────────────────────────────────────
\titleformat{\section}
  {\Large\bfseries\color{blue}}
  {\thesection.}{0.6em}{}
  [\vspace{-4pt}\textcolor{accent}{\rule{\linewidth}{2pt}}\vspace{2pt}]

\titleformat{\subsection}
  {\large\bfseries\color{lblue}}
  {\thesubsection}{0.6em}{}

\titleformat{\subsubsection}
  {\normalsize\bfseries\color{accent}}
  {\thesubsubsection}{0.6em}{}

% ─────────────────────────────────────────────────────────────
% Hyperlinks
% ─────────────────────────────────────────────────────────────
\hypersetup{
  colorlinks=true,
  linkcolor=blue,
  urlcolor=accent,
  citecolor=lblue,
  pdftitle={Rapport S2R - Pipeline Complet},
  pdfauthor={INPT Project Team},
}

% ─────────────────────────────────────────────────────────────
% Metadonnees
% ─────────────────────────────────────────────────────────────
\title{}
\author{}
\date{}

% ═════════════════════════════════════════════════════════════
\begin{document}
% ═════════════════════════════════════════════════════════════

% ─────────────────────────────────────────────────────────────
% PAGE DE COUVERTURE
% ─────────────────────────────────────────────────────────────
\begin{titlepage}
  \pagecolor{navy}
  \begin{center}
    \vspace*{1.8cm}

    \begin{tikzpicture}
      \fill[accent] (0,0) circle (1.2cm);
      \node at (0,0) {\Huge\bfseries\textcolor{white}{S2R}};
    \end{tikzpicture}

    \vspace{1.0cm}

    {\fontsize{26}{32}\selectfont\bfseries\textcolor{white}{Speech-to-Retrieval System}}

    \vspace{0.4cm}

    {\Large\textcolor{mgray}{Pipeline complet de recherche documentaire par la voix}}

    \vspace{0.8cm}

    \textcolor{accent}{\rule{10cm}{1.5pt}}

    \vspace{0.6cm}

    {\large\bfseries\textcolor{accent}{Rapport Technique --- Pipeline P1 $\rightarrow$ P4}}

    \vspace{1.6cm}

    {\large\bfseries\textcolor{white}{P1 Data Engineer \;|\; P2 Speech AI Engineer}}
    \vspace{0.3cm}

    {\large\bfseries\textcolor{white}{P3 ML Engineer \;|\; P4 System Integration}}

    \vspace{0.6cm}

    {\normalsize\textcolor{mgray}{Institut National des Postes et Telecommunications (INPT)}}

    \vspace{0.3cm}

    {\normalsize\textcolor{mgray}{Rabat, Maroc --- 2024--2025}}

    \vspace{1.8cm}

    \begin{tikzpicture}
      \foreach \i/\lbl/\val in {
        0/Dataset/4\,985 audios,
        1/Corpus/6\,744 chunks,
        2/Modele/Wav2Vec2,
        3/Espace/dim = 256
      }{
        \fill[navy!80!black] (\i*3.8, 0) rectangle (\i*3.8+3.4, 1.3);
        \draw[accent, line width=0.6pt] (\i*3.8, 0) rectangle (\i*3.8+3.4, 1.3);
        \node[accent, font=\small\bfseries] at (\i*3.8+1.7, 0.95) {\lbl};
        \node[white, font=\footnotesize] at (\i*3.8+1.7, 0.35) {\val};
      }
    \end{tikzpicture}

    \vfill

    {\small\textcolor{mgray}{Projet Deep Learning --- Rapport genere avec \LaTeX}}
  \end{center}
\end{titlepage}

\pagecolor{white}

% ─────────────────────────────────────────────────────────────
% TABLE DES MATIERES
% ─────────────────────────────────────────────────────────────
\tableofcontents
\newpage

% ═════════════════════════════════════════════════════════════
\section{Introduction}
% ═════════════════════════════════════════════════════════════

Ce rapport presente le travail realise dans le cadre du projet
\textit{Speech-to-Retrieval System (S2R) using Deep Learning}, developpe a l'Institut
National des Postes et Telecommunications (INPT). Il couvre l'integralite des quatre
parties du pipeline :

\begin{itemize}[leftmargin=1.5cm]
  \item \textbf{P1 --- Data Engineer} : preparation du dataset speech $\rightarrow$ document
  \item \textbf{P2 --- Speech AI Engineer} : encodage vocal sans ASR via Wav2Vec2
  \item \textbf{P3 --- ML Engineer} : entrainement du Dual Encoder (alignement audio--texte)
  \item \textbf{P4 --- System Integration} : API REST FastAPI, interface Gradio, deploiement
\end{itemize}

Un systeme Speech-to-Retrieval (S2R) permet de retrouver des documents textuels pertinents
a partir d'une requete vocale, \textbf{sans passer par une etape de transcription ASR}.
L'audio est directement converti en vecteur d'embedding dans un espace semantique partage
avec les documents textuels, ce qui rend le systeme robuste aux accents et aux langues mixtes.

\begin{table}[H]
\centering
\caption{Vue d'ensemble du pipeline S2R}
\begin{tabularx}{\linewidth}{>{\bfseries}l l l l}
\toprule
\rowcolor{navy}\textcolor{white}{Partie} & \textcolor{white}{Module} &
\textcolor{white}{Role} & \textcolor{white}{Statut} \\
\midrule
P1 & \texttt{notebooks/P1\_*.ipynb}        & Donnees audio + texte   & \textcolor{green}{\textbf{\checkmark\ Termine}} \\
\rowcolor{silver}
P2 & \texttt{src/speech\_encoder.py}        & Encodage audio Wav2Vec2 & \textcolor{green}{\textbf{\checkmark\ Termine}} \\
P3 & \texttt{training/}                      & Dual Encoder (InfoNCE)  & \textcolor{green}{\textbf{\checkmark\ Termine}} \\
\rowcolor{silver}
P4 & \texttt{api/ + demo/}                  & API FastAPI + Gradio    & \textcolor{green}{\textbf{\checkmark\ Termine}} \\
\bottomrule
\end{tabularx}
\end{table}

\subsection{Architecture globale}

\begin{figure}[H]
\centering
\begin{tikzpicture}[
  node distance=0.6cm and 1.4cm,
  box/.style={rectangle, draw=blue!60, fill=blue!8, rounded corners=4pt,
              minimum width=4.2cm, minimum height=0.75cm,
              font=\small, align=center},
  gbox/.style={rectangle, draw=green!60, fill=green!8, rounded corners=4pt,
               minimum width=4.2cm, minimum height=0.75cm,
               font=\small, align=center},
  arrow/.style={-{Stealth[length=5pt]}, thick, color=accent}
]
  % Colonne audio
  \node[box, fill=navy!85, text=white] (wav) {.wav brut (16 kHz)};
  \node[box, below=of wav] (w2v) {Wav2Vec2\\(768D, gele)};
  \node[box, below=of w2v] (aproj) {Audio Projection\\(Linear + Norm L2)};
  \node[gbox, below=of aproj] (qemb) {Query Embedding\\(256D)};

  % Colonne texte
  \node[box, right=3.5cm of wav, fill=navy!85, text=white] (txt) {Corpus texte (chunks)};
  \node[box, below=of txt] (mini) {MiniLM\\(384D, gele)};
  \node[box, below=of mini] (tproj) {Text Projection\\(Linear + Norm L2)};
  \node[gbox, below=of tproj] (temb) {Text Embeddings\\(256D)};

  % FAISS
  \node[gbox, fill=accent!20, draw=accent, below=2.0cm of qemb, xshift=2.6cm]
    (faiss) {FAISS IndexFlatIP};

  % Resultats
  \node[gbox, fill=green!15, draw=green!50, below=0.6cm of faiss]
    (res) {Top-k documents};

  \draw[arrow] (wav)   -- (w2v);
  \draw[arrow] (w2v)   -- (aproj);
  \draw[arrow] (aproj) -- (qemb);
  \draw[arrow] (txt)   -- (mini);
  \draw[arrow] (mini)  -- (tproj);
  \draw[arrow] (tproj) -- (temb);
  \draw[arrow] (qemb)  -- (faiss);
  \draw[arrow] (temb)  -- (faiss);
  \draw[arrow] (faiss) -- (res);
\end{tikzpicture}
\caption{Architecture complete du systeme Speech-to-Retrieval}
\label{fig:global}
\end{figure}

\newpage

% ═════════════════════════════════════════════════════════════
\section{P1 --- Data Engineer}
% ═════════════════════════════════════════════════════════════

La partie P1 est responsable de la construction et de la preparation du dataset
\textit{speech $\rightarrow$ document}. Elle fournit l'ensemble des donnees necessaires
aux parties P2 et P3.

% ─────────────────────────────────────────────────────────────
\subsection{Collecte des donnees audio}
% ─────────────────────────────────────────────────────────────

Les donnees audio ont ete collectees depuis \textbf{openslr/librispeech\_asr}
(HuggingFace), un dataset anglophone issu de livres audio du domaine public.
Deux splits ont ete utilises :

\begin{itemize}[leftmargin=1.5cm]
  \item \textbf{\texttt{train.100}} --- 100 heures d'audio propre, 3\,000 echantillons
  \item \textbf{\texttt{train.360}} --- 360 heures d'audio propre, 2\,000 echantillons
\end{itemize}

\begin{table}[H]
\centering
\caption{Caracteristiques du dataset audio brut}
\begin{tabularx}{\linewidth}{>{\bfseries\color{blue}}l X}
\toprule
Parametre & Valeur \\
\midrule
Dataset        & openslr/librispeech\_asr (HuggingFace) \\
\rowcolor{silver}
Langue         & Anglais (English) \\
Splits         & train.100 + train.360 \\
\rowcolor{silver}
Format brut    & Divers sample rates (16 kHz, 22 kHz, 44 kHz) \\
Acces          & Public --- aucune authentification requise \\
\bottomrule
\end{tabularx}
\end{table}

% ─────────────────────────────────────────────────────────────
\subsection{Pretraitement audio}
% ─────────────────────────────────────────────────────────────

Chaque fichier audio a subi un pipeline de nettoyage standardise avant d'etre
sauvegarde en format \textbf{WAV 16\,kHz mono} :

\begin{enumerate}[leftmargin=1.5cm]
  \item \textbf{Resample} --- Conversion vers 16\,000\,Hz via \texttt{librosa.resample()}
  \item \textbf{Mono} --- Conversion stereo $\rightarrow$ mono
  \item \textbf{Filtre duree} --- Rejet des audios $<$ 1\,s, troncature a 20\,s
  \item \textbf{Filtre SNR} --- Rejet des audios avec SNR $<$ 5\,dB
  \item \textbf{Normalisation peak} --- Normalisation amplitude a $[-1, 1]$
\end{enumerate}

\begin{lstlisting}[caption={Estimation du SNR (Signal-to-Noise Ratio)}, label={lst:snr}]
signal_power = np.mean(waveform ** 2) + 1e-10
noise_floor  = np.percentile(np.abs(waveform), 10) ** 2 + 1e-10
snr_db = 10 * np.log10(signal_power / noise_floor)
if snr_db < MIN_SNR_DB:  # seuil = 5 dB
    return None  # audio rejete
\end{lstlisting}

\begin{resultbox}
\textbf{Resultat :} 4\,985 fichiers WAV acceptes sur 5\,000 cibles (99.7\%).
Les 15 fichiers rejetes presentaient un SNR insuffisant ou une duree trop courte.
\end{resultbox}

% ─────────────────────────────────────────────────────────────
\subsection{Construction du corpus texte}
% ─────────────────────────────────────────────────────────────

Le corpus textuel a ete construit a partir d'articles Wikipedia en anglais via
\textbf{wikimedia/wikipedia (20231101.en)} : 600 articles telecharges en streaming,
filtrage des articles $<$ 150 mots, \textbf{570 articles retenus}.

Chaque article a ete decoupe en chunks de \textbf{512 tokens} avec un overlap de
\textbf{50 tokens} (fenetre glissante) :

\begin{lstlisting}[caption={Pipeline de chunking avec fenetre glissante}]
tokenizer = AutoTokenizer.from_pretrained(
    "sentence-transformers/all-mpnet-base-v2"
)
tokens = tokenizer.encode(document_text, add_special_tokens=False)
chunk_size, overlap = 512, 50
for start in range(0, len(tokens), chunk_size - overlap):
    end   = min(start + chunk_size, len(tokens))
    chunk = tokenizer.decode(tokens[start:end])
\end{lstlisting}

\begin{table}[H]
\centering
\caption{Statistiques du corpus texte}
\begin{tabular}{lrr}
\toprule
\textbf{Metrique} & \textbf{Valeur} & \textbf{Unite} \\
\midrule
Total chunks produits  & 6\,744  & chunks \\
\rowcolor{silver}
Articles source        & 570     & documents \\
Tokens moyens/chunk    & 493     & tokens \\
\rowcolor{silver}
Taille fichier CSV     & 15\,719 & KB \\
\bottomrule
\end{tabular}
\end{table}

% ─────────────────────────────────────────────────────────────
\subsection{Creation des paires audio--document et split}
% ─────────────────────────────────────────────────────────────

Chaque fichier audio est associe au chunk texte le plus pertinent via la
\textbf{similarite de Jaccard} sur les mots-cles :

\begin{equation}
  J(A, B) = \frac{|KW(A) \cap KW(B)|}{|KW(A) \cup KW(B)|}
\end{equation}

\begin{infobox}[Score Jaccard moyen : 0.073]
Score faible attendu : les transcriptions LibriSpeech sont des phrases courtes
tandis que les chunks Wikipedia sont des paragraphes techniques denses.
Les paires servent uniquement a l'entrainement contrastif de P3.
\end{infobox}

\begin{table}[H]
\centering
\caption{Repartition train / val / test (strategie 80/10/10)}
\begin{tabular}{lrrr}
\toprule
\rowcolor{navy}
\textcolor{white}{\textbf{Split}} &
\textcolor{white}{\textbf{Paires}} &
\textcolor{white}{\textbf{Proportion}} &
\textcolor{white}{\textbf{Fichier}} \\
\midrule
Train      & 3\,988 & 80\% & \texttt{pairs\_train.csv} \\
\rowcolor{silver}
Validation & 498    & 10\% & \texttt{pairs\_val.csv} \\
Test       & 499    & 10\% & \texttt{pairs\_test.csv} \\
\midrule
\rowcolor{silver}
\textbf{Total} & \textbf{4\,985} & \textbf{100\%} & \texttt{pairs.csv} \\
\bottomrule
\end{tabular}
\end{table}

\newpage

% ═════════════════════════════════════════════════════════════
\section{P2 --- Speech AI Engineer}
% ═════════════════════════════════════════════════════════════

La partie P2 transforme les fichiers audio en vecteurs d'embeddings denses de dimension
\textbf{768}, normalises L2, directement exploitables par le Dual Encoder de P3.

% ─────────────────────────────────────────────────────────────
\subsection{Modele Wav2Vec2}
% ─────────────────────────────────────────────────────────────

Le modele utilise est \textbf{facebook/wav2vec2-base-960h}, pre-entraine par Meta AI
sur 960 heures d'audio anglais (LibriSpeech). C'est un Transformer auto-supervise
qui apprend des representations speech riches sans supervision explicite.

\begin{table}[H]
\centering
\caption{Parametres du modele Wav2Vec2}
\begin{tabular}{>{\bfseries\color{blue}}l l l}
\toprule
Parametre & Valeur & Description \\
\midrule
Modele       & wav2vec2-base-960h & Version base \\
\rowcolor{silver}
Parametres   & 94M               & Total parametres \\
Sample rate  & 16\,000 Hz        & Format audio requis \\
\rowcolor{silver}
Hidden dim   & 768               & Dimension des hidden states \\
Nb couches   & 12                & Transformer layers \\
\rowcolor{silver}
Nb tetes     & 12                & Attention heads \\
Poids        & Geles             & \texttt{frozen=True} pendant training \\
\rowcolor{silver}
Pooling      & Mean              & Moyenne sur dimension temporelle \\
\bottomrule
\end{tabular}
\end{table}

% ─────────────────────────────────────────────────────────────
\subsection{Architecture du SpeechEncoder}
% ─────────────────────────────────────────────────────────────

\begin{figure}[H]
\centering
\begin{tikzpicture}[
  node distance=1.0cm,
  box/.style={rectangle, draw=blue!60, fill=blue!8, rounded corners=4pt,
              minimum width=7cm, minimum height=0.75cm,
              font=\small, align=center},
  arrow/.style={-{Stealth[length=6pt]}, thick, color=accent}
]
  \node[box, fill=navy!90, text=white, font=\small\bfseries] (in)
    {Fichier .wav (16 kHz, mono, float32)};
  \node[box, below=of in] (step1)
    {Chargement via soundfile (sans FFmpeg)};
  \node[box, below=of step1] (step2)
    {Feature extraction --- Wav2Vec2Processor};
  \node[box, below=of step2] (step3)
    {Encodage --- Wav2Vec2Model (12 layers)};
  \node[box, below=of step3, font=\small\itshape] (step3b)
    {hidden states $\in \mathbb{R}^{B \times T' \times 768}$};
  \node[box, below=of step3b] (step4)
    {Mean pooling temporel (avec masque optionnel)};
  \node[box, fill=green!15, draw=green!50, below=of step4, font=\small\bfseries] (out)
    {Embedding norme L2 $\in \mathbb{R}^{768}$};
  \draw[arrow] (in) -- (step1);
  \draw[arrow] (step1) -- (step2);
  \draw[arrow] (step2) -- (step3);
  \draw[arrow] (step3) -- (step3b);
  \draw[arrow] (step3b) -- (step4);
  \draw[arrow] (step4) -- (out);
\end{tikzpicture}
\caption{Pipeline d'encodage audio P2}
\label{fig:p2pipeline}
\end{figure}

% ─────────────────────────────────────────────────────────────
\subsection{Gestion du chargement audio (soundfile)}
% ─────────────────────────────────────────────────────────────

Sur Windows, la dependance \texttt{torchaudio.load} requiert FFmpeg via
\texttt{torchcodec}, ce qui genere des erreurs sur certains environnements.
La solution adoptee utilise directement \textbf{soundfile} :

\begin{lstlisting}[caption={Chargement audio robuste (src/speech\_encoder.py)}]
import soundfile as sf

def _load_wav(audio_path, target_sr=16_000):
    data, sr = sf.read(audio_path, dtype="float32", always_2d=True)
    waveform = torch.from_numpy(data.T)  # (channels, samples)
    if waveform.shape[0] > 1:            # stereo -> mono
        waveform = waveform.mean(dim=0, keepdim=True)
    if sr != target_sr:
        # torchaudio.functional.resample = op PyTorch pur, sans backend
        waveform = torchaudio.functional.resample(waveform, sr, target_sr)
    return waveform.squeeze(0)           # [T]
\end{lstlisting}

% ─────────────────────────────────────────────────────────────
\subsection{Mean pooling avec masque d'attention}
% ─────────────────────────────────────────────────────────────

Pour les batches paddes, un masque d'attention est propage jusqu'a la sortie
du feature extractor pour un pooling correct :

\begin{equation}
  \mathbf{e} = \frac{\sum_{t=1}^{T'} \mathbf{h}_t \cdot m_t}{\sum_{t=1}^{T'} m_t}
  \qquad m_t \in \{0, 1\}
\end{equation}

La normalisation L2 finale garantit $\|\mathbf{e}\|_2 = 1$, rendant la similarite
cosine equivalente au produit scalaire, ce qui est optimal pour FAISS IndexFlatIP.

\begin{equation}
  \hat{\mathbf{e}} = \frac{\mathbf{e}}{\|\mathbf{e}\|_2}
  \implies \cos(\hat{\mathbf{e}}_1, \hat{\mathbf{e}}_2) = \hat{\mathbf{e}}_1^\top \hat{\mathbf{e}}_2
\end{equation}

% ─────────────────────────────────────────────────────────────
\subsection{Outputs de P2}
% ─────────────────────────────────────────────────────────────

\begin{table}[H]
\centering
\caption{Fichiers produits par P2}
\begin{tabularx}{\linewidth}{>{\ttfamily}l l X}
\toprule
\textbf{Fichier} & \textbf{Dimensions} & \textbf{Description} \\
\midrule
audio\_embeddings.npy        & $(4985, 768)$ & Matrice des embeddings Wav2Vec2 \\
\rowcolor{silver}
audio\_embeddings\_index.csv  & 4985 lignes  & Mapping audio\_id $\rightarrow$ ligne \\
src/speech\_encoder.py        & ---          & Module importable par P3/P4 \\
\bottomrule
\end{tabularx}
\end{table}

\newpage

% ═════════════════════════════════════════════════════════════
\section{P3 --- ML Engineer : Dual Encoder}
% ═════════════════════════════════════════════════════════════

La partie P3 entraine un \textbf{Dual Encoder} qui aligne l'espace audio (768D) et
l'espace texte (384D) dans un espace commun partage de dimension \textbf{256D},
via une fonction de perte contrastive.

% ─────────────────────────────────────────────────────────────
\subsection{Architecture du DualEncoderModel}
% ─────────────────────────────────────────────────────────────

Le modele est compose de deux branches de projection symetriques :

\begin{figure}[H]
\centering
\begin{tikzpicture}[
  node distance=0.8cm and 2.0cm,
  box/.style={rectangle, draw=blue!60, fill=blue!8, rounded corners=4pt,
              minimum width=3.8cm, minimum height=0.7cm,
              font=\small, align=center},
  gbox/.style={rectangle, draw=green!60, fill=green!8, rounded corners=4pt,
               minimum width=3.8cm, minimum height=0.7cm,
               font=\small, align=center},
  arrow/.style={-{Stealth[length=5pt]}, thick, color=accent}
]
  % Branche audio
  \node[box, fill=navy!85, text=white] (a0) {Wav2Vec2 embedding\\$\mathbb{R}^{768}$ (gele)};
  \node[box, below=of a0] (a1) {Dropout($p=0.1$)};
  \node[box, below=of a1] (a2) {Linear($768 \to 256$)};
  \node[gbox, below=of a2] (az) {$z_s \in \mathbb{R}^{256}$, L2-norme};

  % Branche texte
  \node[box, fill=navy!85, text=white, right=2.5cm of a0] (t0) {MiniLM embedding\\$\mathbb{R}^{384}$ (pre-calcule)};
  \node[box, below=of t0] (t1) {Dropout($p=0.1$)};
  \node[box, below=of t1] (t2) {Linear($384 \to 256$)};
  \node[gbox, below=of t2] (tz) {$z_t \in \mathbb{R}^{256}$, L2-norme};

  % Similarite
  \node[gbox, fill=accent!20, draw=accent, below=1.2cm of a2, xshift=2.6cm]
    (sim) {$\text{sim}(z_s, z_t) = z_s^\top z_t$};

  \draw[arrow] (a0) -- (a1);
  \draw[arrow] (a1) -- (a2);
  \draw[arrow] (a2) -- (az);
  \draw[arrow] (t0) -- (t1);
  \draw[arrow] (t1) -- (t2);
  \draw[arrow] (t2) -- (tz);
  \draw[arrow] (az)  -- (sim);
  \draw[arrow] (tz)  -- (sim);
\end{tikzpicture}
\caption{Architecture du DualEncoderModel (mode rapide)}
\label{fig:dualenc}
\end{figure}

\begin{lstlisting}[caption={DualEncoderModel simplifie (training/models.py)}]
class DualEncoderModel(nn.Module):
    def __init__(self, audio_input_dim=768, projection_dim=256,
                 dropout=0.1, text_input_dim=384):
        super().__init__()
        self.text_encoder = None  # mode rapide
        self.audio_projection = nn.Sequential(
            nn.Dropout(dropout),
            nn.Linear(audio_input_dim, projection_dim),
        )
        self.text_projection = nn.Sequential(
            nn.Dropout(dropout),
            nn.Linear(text_input_dim, projection_dim),
        )

    def encode_audio(self, x):
        return F.normalize(self.audio_projection(x), p=2, dim=-1)

    def encode_text_precomputed(self, x):
        return F.normalize(self.text_projection(x), p=2, dim=-1)
\end{lstlisting}

% ─────────────────────────────────────────────────────────────
\subsection{Mode rapide (pre-calcul des embeddings texte)}
% ─────────────────────────────────────────────────────────────

Sur CPU, encoder MiniLM a chaque batch pendant l'entrainement serait prohibitif.
La solution adoptee est le \textbf{mode rapide} :

\begin{enumerate}[leftmargin=1.5cm]
  \item Pre-calculer \textit{une seule fois} tous les embeddings MiniLM (384D) en
        \texttt{embeddings/precomputed\_text.npy}
  \item Pendant l'entrainement, charger ces vecteurs depuis la memoire (numpy array)
  \item Uniquement les deux couches \texttt{Linear} sont entrainables
\end{enumerate}

\begin{table}[H]
\centering
\caption{Comparaison mode complet vs mode rapide}
\begin{tabular}{lll}
\toprule
\textbf{Critere} & \textbf{Mode complet} & \textbf{Mode rapide} \\
\midrule
Forward text/batch & MiniLM complet & Lookup numpy \\
\rowcolor{silver}
Temps/epoque (CPU) & $\sim$30 min    & $\sim$45 secondes \\
Parametres entrain. & 22M (MPNet)   & \textbf{2 couches Linear} \\
\rowcolor{silver}
Requiert GPU       & Oui            & Non \\
Speedup            & 1$\times$      & \textbf{20--50$\times$} \\
\bottomrule
\end{tabular}
\end{table}

\begin{resultbox}
\textbf{Acceleration mesuree sur Intel Core i7, 24 Go RAM :}
mode complet $\approx$ 2h20/epoque $\rightarrow$ mode rapide $\approx$ 45 sec/epoque.
Entrainement complet de 20 epoques ramene de $\approx$ 2 jours a $\approx$ 15 minutes.
\end{resultbox}

% ─────────────────────────────────────────────────────────────
\subsection{Fonction de perte}
% ─────────────────────────────────────────────────────────────

Deux fonctions de perte sont implementees dans \texttt{training/losses.py}.

\subsubsection{Contrastive Loss (InfoNCE / NT-Xent)}

Pour un batch de $N$ paires $(z_s^{(i)}, z_t^{(i)})$, la matrice de similarites est :

\begin{equation}
  S_{ij} = \frac{z_s^{(i)} \cdot z_t^{(j)}}{\tau}, \quad \tau = 0.07
\end{equation}

La perte InfoNCE penalise les similarites hors-diagonale (negatifs in-batch) :

\begin{equation}
  \mathcal{L} =
  \frac{1}{2}
  \left(
    \underbrace{\frac{1}{N}\sum_{i=1}^{N} -\log \frac{e^{S_{ii}}}{\sum_j e^{S_{ij}}}}_{\mathcal{L}_{s \to t}}
    +
    \underbrace{\frac{1}{N}\sum_{i=1}^{N} -\log \frac{e^{S_{ii}}}{\sum_j e^{S_{ji}}}}_{\mathcal{L}_{t \to s}}
  \right)
\end{equation}

\begin{lstlisting}[caption={Contrastive loss (InfoNCE) -- training/losses.py}]
def contrastive_loss(speech_emb, text_emb, temperature=0.07):
    logits = (speech_emb @ text_emb.T) / temperature
    labels = torch.arange(logits.size(0), device=logits.device)
    loss_s2t = F.cross_entropy(logits,   labels)
    loss_t2s = F.cross_entropy(logits.T, labels)
    return (loss_s2t + loss_t2s) / 2.0
\end{lstlisting}

\subsubsection{Triplet Loss}

Comme alternative, une triplet loss avec hard-negative mining est disponible :

\begin{equation}
  \mathcal{L}_{triplet} = \frac{1}{2}\Big(
  \text{ReLU}\!\left(d^+ - d^-_{text} + m\right)
  + \text{ReLU}\!\left(d^+ - d^-_{speech} + m\right)
  \Big)
\end{equation}

avec $d^+ = 1 - \text{sim}(z_s^{(i)}, z_t^{(i)})$ et $m = 0.2$.

% ─────────────────────────────────────────────────────────────
\subsection{Procedure d'entrainement}
% ─────────────────────────────────────────────────────────────

\begin{table}[H]
\centering
\caption{Hyperparametres d'entrainement recommandes}
\begin{tabular}{>{\bfseries}l l l}
\toprule
Parametre & Valeur locale (CPU) & Valeur cloud (GPU T4) \\
\midrule
Epoques          & 20      & 20 \\
\rowcolor{silver}
Batch size       & 16      & 64 \\
Learning rate    & 2e-4    & 2e-4 \\
\rowcolor{silver}
Warmup steps     & 100     & 100 \\
Grad clipping    & 1.0     & 1.0 \\
\rowcolor{silver}
Projection dim   & 256     & 256 \\
Temperature      & 0.07    & 0.07 \\
\rowcolor{silver}
Optimiseur       & AdamW   & AdamW \\
\bottomrule
\end{tabular}
\end{table}

Le scheduler de learning rate combine un warmup lineaire sur les premiers 100 steps
puis un decroissance cosine jusqu'a la fin de l'entrainement.

% ─────────────────────────────────────────────────────────────
\subsection{Metriques d'evaluation}
% ─────────────────────────────────────────────────────────────

Trois metriques sont calculees a chaque epoque de validation
dans \texttt{training/metrics.py} :

\begin{itemize}[leftmargin=1.5cm]
  \item \textbf{Recall@k} : proportion de requetes pour lesquelles le document correct
        apparait dans le top-$k$ (k = 5 et 10)

  \begin{equation}
    \text{Recall@}k = \frac{1}{|Q|}\sum_{q \in Q} \mathbf{1}\left[\text{rank}(q) \leq k\right]
  \end{equation}

  \item \textbf{MRR} (Mean Reciprocal Rank) : moyenne des inverses du rang du premier
        document pertinent

  \begin{equation}
    \text{MRR} = \frac{1}{|Q|}\sum_{q \in Q} \frac{1}{\text{rank}(q)}
  \end{equation}
\end{itemize}

\begin{table}[H]
\centering
\caption{Objectifs de performance apres 20 epoques}
\begin{tabular}{lll}
\toprule
\textbf{Metrique} & \textbf{Objectif} & \textbf{Interpretation} \\
\midrule
Recall@5  & $> 0.30$ & 30\% des requetes trouvees dans le top-5 \\
\rowcolor{silver}
Recall@10 & $> 0.45$ & 45\% des requetes trouvees dans le top-10 \\
MRR       & $> 0.25$ & rang moyen reciproque du document correct \\
\bottomrule
\end{tabular}
\end{table}

\begin{infobox}[Limitation intrinseque du dataset]
Les paires audio--document sont construites par similarite Jaccard (score moy. 0.073).
Les paires ne sont pas semantiquement parfaites : les metriques mesurent la capacite
du modele a retrouver ces paires imparfaites, pas la qualite semantique absolue.
\end{infobox}

% ─────────────────────────────────────────────────────────────
\subsection{Export des embeddings texte (post-entrainement)}
% ─────────────────────────────────────────────────────────────

Apres l'entrainement, \texttt{training/export\_text\_embeddings.py} projette
tous les chunks texte dans l'espace commun 256D :

\begin{enumerate}[leftmargin=1.5cm]
  \item Charger les embeddings pre-calcules (384D) depuis \texttt{precomputed\_text.npy}
  \item Appliquer \texttt{model.encode\_text\_precomputed()} par batch de 64
  \item Sauvegarder la matrice $(N, 256)$ dans \texttt{text\_chunk\_embeddings.npy}
  \item Sauvegarder le manifest (\texttt{chunk\_id}, \texttt{document\_id}, \texttt{text})
\end{enumerate}

\newpage

% ═════════════════════════════════════════════════════════════
\section{P4 --- System Integration}
% ═════════════════════════════════════════════════════════════

La partie P4 expose le systeme S2R via deux interfaces :
une \textbf{API REST FastAPI} pour l'integration backend,
et une \textbf{interface Gradio} pour la demonstration interactive.

% ─────────────────────────────────────────────────────────────
\subsection{Index FAISS}
% ─────────────────────────────────────────────────────────────

La recherche des documents pertinents est effectuee via
\textbf{FAISS IndexFlatIP} (produit scalaire exact sur vecteurs normalises = cosine similarity).
L'index est construit en memoire au demarrage :

\begin{lstlisting}[caption={Construction de l'index FAISS (api/main.py)}]
embs = np.load("embeddings/text_chunk_embeddings.npy").astype("float32")
norms = np.linalg.norm(embs, axis=1, keepdims=True)
embs /= np.clip(norms, 1e-9, None)          # normalisation L2

index = faiss.IndexFlatIP(embs.shape[1])     # dim = 256
index.add(embs)                              # N vecteurs indexes
# Recherche : scores, indices = index.search(query, k)
\end{lstlisting}

% ─────────────────────────────────────────────────────────────
\subsection{API REST FastAPI}
% ─────────────────────────────────────────────────────────────

L'API expose 5 endpoints documentes via Swagger (disponible sur \texttt{/docs}) :

\begin{table}[H]
\centering
\caption{Endpoints de l'API FastAPI (api/main.py)}
\begin{tabularx}{\linewidth}{l l X}
\toprule
\textbf{Methode} & \textbf{Route} & \textbf{Description} \\
\midrule
GET  & \texttt{/health}        & Etat du service (dual encoder charge, taille index) \\
\rowcolor{silver}
POST & \texttt{/search/audio}  & Upload fichier .wav $\to$ top-k documents \\
POST & \texttt{/search/base64} & Audio encode base64 $\to$ top-k documents \\
\rowcolor{silver}
POST & \texttt{/search/text}   & Requete texte (MiniLM) $\to$ top-k documents \\
GET  & \texttt{/examples}      & Liste des fichiers audio de demonstration \\
\bottomrule
\end{tabularx}
\end{table}

\begin{lstlisting}[caption={Exemple de reponse /search/audio}]
{
  "results": [
    {"rank": 1, "score": 0.8241, "chunk_id": "doc00042_chunk003",
     "text": "Anarchism is a political philosophy..."},
    {"rank": 2, "score": 0.7955, "chunk_id": "doc00042_chunk004",
     "text": "The anarchist movement emerged..."}
  ],
  "mode": "dual_encoder",
  "query_dim": 256,
  "index_size": 6744
}
\end{lstlisting}

Lancement de l'API :

\begin{lstlisting}[language=bash, caption={Lancement de l'API FastAPI}]
# Variables d'environnement (optionnelles)
export S2R_CHECKPOINT="models/dual_encoder_mpnet/best_model.pt"
export S2R_TEXT_EMBEDDINGS="embeddings/text_chunk_embeddings.npy"
export S2R_MANIFEST="embeddings/text_chunk_manifest.csv"

# Demarrage
uvicorn api.main:app --host 0.0.0.0 --port 8000 --reload
# Documentation interactive : http://localhost:8000/docs
\end{lstlisting}

% ─────────────────────────────────────────────────────────────
\subsection{Interface Gradio}
% ─────────────────────────────────────────────────────────────

L'interface Gradio (\texttt{demo/app.py}) propose deux onglets :

\begin{itemize}[leftmargin=1.5cm]
  \item \textbf{Onglet 1 --- Recherche vocale} : upload ou enregistrement microphone,
        slider pour le nombre de resultats, liste deroulante d'exemples pre-charges
  \item \textbf{Onglet 2 --- Recherche texte} : saisie textuelle directe,
        comparaison avec le mode audio pour valider l'index
\end{itemize}

\begin{table}[H]
\centering
\caption{Options de lancement de la demo Gradio}
\begin{tabular}{l l}
\toprule
\textbf{Commande} & \textbf{Description} \\
\midrule
\texttt{python demo/app.py} & Mode prototype (Wav2Vec2 brut) \\
\rowcolor{silver}
\texttt{python demo/app.py --checkpoint best\_model.pt} & Dual encoder entraine \\
\texttt{python demo/app.py --share} & Lien public Gradio (72h) \\
\rowcolor{silver}
\texttt{python demo/app.py --port 7861} & Port personnalise \\
\bottomrule
\end{tabular}
\end{table}

Au demarrage, la demo affiche un panneau de statut indiquant le mode actif
(prototype ou dual encoder entraine), la taille de l'index FAISS, et le device utilise.

\newpage

% ═════════════════════════════════════════════════════════════
\section{Pipeline d'execution}
% ═════════════════════════════════════════════════════════════

% ─────────────────────────────────────────────────────────────
\subsection{Pipeline automatise (run\_pipeline.bat / .sh)}
% ─────────────────────────────────────────────────────────────

Deux scripts d'automatisation complets sont fournis pour Windows et Linux/macOS.
Le pipeline s'execute en 6 etapes sequentielles :

\begin{table}[H]
\centering
\caption{Etapes du pipeline automatise}
\begin{tabularx}{\linewidth}{c l X}
\toprule
\textbf{Etape} & \textbf{Module} & \textbf{Action} \\
\midrule
1 & \texttt{validate\_handoff} & Verification des donnees P1/P2 \\
\rowcolor{silver}
2 & \texttt{precompute\_text}  & Pre-calcul embeddings MiniLM (une fois) \\
3 & \texttt{train\_dual\_encoder} & Entrainement 20 epoques, batch 16 \\
\rowcolor{silver}
4 & \texttt{export\_text\_embeddings} & Export embeddings 256D vers FAISS \\
5 & \texttt{inference.py}     & Test d'inference sur un fichier exemple \\
\rowcolor{silver}
6 & API ou Demo (optionnel)   & Lancement selon \texttt{--api} ou \texttt{--demo} \\
\bottomrule
\end{tabularx}
\end{table}

\begin{lstlisting}[language=bash, caption={Options du pipeline Windows}]
run_pipeline.bat                    # Pipeline complet (20 epoques)
run_pipeline.bat --skip-train       # Sauter l'entrainement (checkpoint present)
run_pipeline.bat --skip-train --demo  # Sauter + lancer Gradio
run_pipeline.bat --skip-train --api   # Sauter + lancer FastAPI
\end{lstlisting}

% ─────────────────────────────────────────────────────────────
\subsection{Entrainement sur Google Colab / Kaggle}
% ─────────────────────────────────────────────────────────────

Un notebook Jupyter complet est fourni dans \texttt{notebooks/S2R\_Colab\_Kaggle.ipynb}
pour l'entrainement sur GPU cloud :

\begin{table}[H]
\centering
\caption{Performances d'entrainement selon la plateforme}
\begin{tabular}{llll}
\toprule
\textbf{Plateforme} & \textbf{GPU} & \textbf{Temps/epoque} & \textbf{20 epoques} \\
\midrule
Local (Intel i7, CPU) & Aucun    & $\approx$ 45 s & $\approx$ 15 min \\
\rowcolor{silver}
Google Colab (gratuit) & T4 15GB & $\approx$ 60 s & $\approx$ 20 min \\
Kaggle (gratuit)       & T4 16GB & $\approx$ 55 s & $\approx$ 18 min \\
\rowcolor{silver}
Google Colab Pro       & A100 40GB & $\approx$ 20 s & $\approx$ 7 min \\
\bottomrule
\end{tabular}
\end{table}

Le notebook couvre : installation des dependances, chargement depuis Google Drive
ou upload direct, validation des donnees, pre-calcul MiniLM, entrainement avec
tqdm progress, evaluation, et telechargement du checkpoint.

% ─────────────────────────────────────────────────────────────
\subsection{Verification des metriques}
% ─────────────────────────────────────────────────────────────

\begin{lstlisting}[caption={Validation des donnees P1/P2 (training/validate\_handoff.py)}]
python -m training.validate_handoff \
    --audio-manifest-csv  embeddings/audio_embeddings_index.csv \
    --audio-embeddings-npy embeddings/audio_embeddings.npy \
    --pairs-csv           data/output/pairs_train.csv

# Attendu :
# manifest_rows=4985
# embedding_shape=(4985, 768)
# status=manifest_and_embeddings_ok
# pairs_rows=3988
# status=pairs_alignment_ok
\end{lstlisting}

\newpage

% ═════════════════════════════════════════════════════════════
\section{Integration Snowflake}
% ═════════════════════════════════════════════════════════════

% ─────────────────────────────────────────────────────────────
\subsection{Architecture SQL}
% ─────────────────────────────────────────────────────────────

Un pipeline SQL complet pour Snowflake est fourni dans \texttt{snowflake\_pipeline.sql}.
Il exploite le type natif \textbf{VECTOR(FLOAT, 256)} de Snowflake
et la fonction \textbf{VECTOR\_COSINE\_SIMILARITY} pour la recherche.

\begin{table}[H]
\centering
\caption{Tables principales du pipeline Snowflake}
\begin{tabularx}{\linewidth}{>{\ttfamily}l X}
\toprule
\textbf{Table} & \textbf{Description} \\
\midrule
AUDIO\_EMBEDDINGS   & Embeddings audio 256D + metadonnees (\texttt{audio\_id}, \texttt{filepath}) \\
\rowcolor{silver}
TEXT\_EMBEDDINGS    & Embeddings texte 256D + contenu des chunks \\
AUDIO\_QUERIES      & Requetes audio de test \\
\rowcolor{silver}
RETRIEVAL\_RESULTS  & Resultats de recherche avec scores et rangs \\
\bottomrule
\end{tabularx}
\end{table}

\begin{lstlisting}[language=SQL, caption={Schema de la table TEXT\_EMBEDDINGS dans Snowflake}]
CREATE TABLE S2R_DB.S2R_SCHEMA.TEXT_EMBEDDINGS (
    chunk_id      VARCHAR(255) NOT NULL PRIMARY KEY,
    document_id   VARCHAR(255),
    text_content  TEXT,
    embedding     VECTOR(FLOAT, 256),      -- type natif Snowflake
    loaded_at     TIMESTAMP DEFAULT CURRENT_TIMESTAMP()
);
\end{lstlisting}

% ─────────────────────────────────────────────────────────────
\subsection{Recherche par similarite vectorielle}
% ─────────────────────────────────────────────────────────────

\begin{lstlisting}[language=SQL, caption={Recherche top-5 dans Snowflake}]
SELECT
    t.chunk_id,
    t.document_id,
    LEFT(t.text_content, 300) AS text_preview,
    VECTOR_COSINE_SIMILARITY(t.embedding, q.embedding) AS score
FROM S2R_DB.S2R_SCHEMA.TEXT_EMBEDDINGS t
CROSS JOIN (
    SELECT embedding FROM S2R_DB.S2R_SCHEMA.AUDIO_QUERIES
    WHERE audio_id = 'spontaneous-speech-fr-20300'
) q
ORDER BY score DESC
LIMIT 5;
\end{lstlisting}

% ─────────────────────────────────────────────────────────────
\subsection{Procedure stockee d'inference}
% ─────────────────────────────────────────────────────────────

Une procedure stockee Snowpark integre le modele Python directement dans Snowflake :

\begin{lstlisting}[language=SQL, caption={Procedure d'inference S2R dans Snowflake}]
CREATE OR REPLACE PROCEDURE S2R_DB.S2R_SCHEMA.S2R_SEARCH(
    AUDIO_PATH VARCHAR,
    K INTEGER
)
RETURNS TABLE(chunk_id VARCHAR, score FLOAT, text_preview VARCHAR)
LANGUAGE PYTHON
RUNTIME_VERSION = '3.10'
PACKAGES = ('torch', 'transformers', 'soundfile', 'numpy', 'faiss-cpu')
HANDLER = 's2r_search'
AS $$
def s2r_search(session, audio_path, k):
    import numpy as np
    from src.speech_encoder import speech_to_embedding
    emb = speech_to_embedding(audio_path)
    # ... recherche dans la table TEXT_EMBEDDINGS
$$;
\end{lstlisting}

\newpage

% ═════════════════════════════════════════════════════════════
\section{Problemes techniques et solutions}
% ═════════════════════════════════════════════════════════════

Cette section documente les problemes rencontres pendant le developpement et
les corrections apportees.

\begin{table}[H]
\centering
\caption{Problemes techniques et corrections}
\begin{tabularx}{\linewidth}{p{4.5cm} p{3.5cm} X}
\toprule
\textbf{Probleme} & \textbf{Cause} & \textbf{Solution} \\
\midrule
\texttt{RuntimeError: libtorchcodec} &
FFmpeg manquant sur Windows &
Remplacement de \texttt{torchaudio.load} par \texttt{soundfile.read} \\

\rowcolor{silver}
\texttt{unrecognized arguments: --text-chunks-csv} &
Argument absent de \texttt{validate\_handoff.py} &
Suppression de l'argument des appels au script \\

\texttt{audio\_id column not found} &
Colonne nommee \texttt{audio\_file} dans certains CSVs &
Normalisation automatique \texttt{audio\_file} $\to$ \texttt{audio\_id} \\

\rowcolor{silver}
\texttt{missing column: document\_id} &
Colonne absent de \texttt{corpus\_chunks.csv} &
Derivation automatique depuis \texttt{chunk\_id} (regex) \\

Segfault DataLoader (Windows) &
Multiprocessing incompatible &
\texttt{num\_workers=0, pin\_memory=False} \\

\rowcolor{silver}
.bat non reconnu &
Caracteres accentues (UTF-8) dans \texttt{cmd.exe} &
Reecriture du .bat en ASCII pur, suppression de \texttt{chcp} \\

Temps d'entrainement $>$ 2h/epoque &
Forward MiniLM a chaque batch &
Mode rapide : pre-calcul unique des embeddings texte \\

\rowcolor{silver}
\texttt{DualEncoderModel(**config)} crash &
\texttt{text\_input\_dim} absent du config sauvegarde &
Correction de \texttt{export\_config()} pour sauvegarder \texttt{text\_input\_dim} \\

\bottomrule
\end{tabularx}
\end{table}

\newpage

% ═════════════════════════════════════════════════════════════
\section{Resultats et evaluation}
% ═════════════════════════════════════════════════════════════

\subsection{Bilan quantitatif global}

\begin{table}[H]
\centering
\caption{Bilan quantitatif du projet complet}
\begin{tabular}{lrl}
\toprule
\rowcolor{navy}
\textcolor{white}{\textbf{Metrique}} &
\textcolor{white}{\textbf{Valeur}} &
\textcolor{white}{\textbf{Detail}} \\
\midrule
Fichiers audio WAV       & 4\,985   & 16 kHz, mono, normalises \\
\rowcolor{silver}
Duree totale audio       & $\approx$17.6 h & $\approx$ 63\,600 secondes \\
Chunks texte             & 6\,744   & 512 tokens, overlap 50 \\
\rowcolor{silver}
Paires audio--document   & 4\,985   & Jaccard moyen = 0.073 \\
Paires d'entrainement    & 3\,988   & 80\% du dataset \\
\rowcolor{silver}
Dimension audio (Wav2Vec2) & 768D   & 94M parametres, gele \\
Dimension texte (MiniLM)   & 384D   & 22M parametres, gele \\
\rowcolor{silver}
Espace partage           & 256D     & Dual Encoder projections \\
Parametres entrainables  & $\sim$300K & 2 couches Linear uniquement \\
\rowcolor{silver}
Temps entrainement CPU   & $\sim$15 min & 20 epoques, mode rapide \\
Recall@5 (objectif)      & $>$ 0.30 & Apres 20 epoques \\
\rowcolor{silver}
Recall@10 (objectif)     & $>$ 0.45 & Apres 20 epoques \\
MRR (objectif)           & $>$ 0.25 & Apres 20 epoques \\
\bottomrule
\end{tabular}
\end{table}

\subsection{Tests systeme}

\begin{table}[H]
\centering
\caption{Resultats des tests systeme}
\begin{tabularx}{\linewidth}{X c l}
\toprule
\textbf{Test} & \textbf{Resultat} & \textbf{Critere} \\
\midrule
Output shape SpeechEncoder = $(768,)$  & \textcolor{green}{\checkmark} & \texttt{emb.shape == (768,)} \\
\rowcolor{silver}
Norme L2 $\approx$ 1.0                 & \textcolor{green}{\checkmark} & $\|\hat{e}\|_2 = 1.0000$ \\
Determinisme du SpeechEncoder          & \textcolor{green}{\checkmark} & meme input $\to$ meme output \\
\rowcolor{silver}
Batch encoding $(N, 768)$              & \textcolor{green}{\checkmark} & shape $[4985, 768]$ \\
Validation handoff P1/P2               & \textcolor{green}{\checkmark} & manifest + paires OK \\
\rowcolor{silver}
Pre-calcul MiniLM $(N, 384)$           & \textcolor{green}{\checkmark} & shape $(6744, 384)$ \\
Export embeddings $(N, 256)$           & \textcolor{green}{\checkmark} & shape $(6744, 256)$ \\
\rowcolor{silver}
API /health                            & \textcolor{green}{\checkmark} & \texttt{"status":"ok"} \\
API /search/audio                      & \textcolor{green}{\checkmark} & 5 resultats avec scores \\
\rowcolor{silver}
Interface Gradio                       & \textcolor{green}{\checkmark} & Demarrage + 2 onglets OK \\
\bottomrule
\end{tabularx}
\end{table}

\newpage

% ═════════════════════════════════════════════════════════════
\section{Conclusion et perspectives}
% ═════════════════════════════════════════════════════════════

\subsection{Bilan}

Ce projet realise un systeme complet de \textbf{Speech-to-Retrieval} :
depuis la collecte et le nettoyage des donnees audio et texte (P1),
l'encodage vocal sans ASR via Wav2Vec2 (P2),
l'alignement des espaces multimodaux par Dual Encoder contrastif (P3),
jusqu'a l'exposition via API REST et interface web interactive (P4).

L'innovation principale est le \textbf{mode rapide d'entrainement} : en pre-calculant
les embeddings texte une seule fois, le pipeline rend l'entrainement accessible sur
CPU ordinaire ($\sim$15 minutes pour 20 epoques), ce qui facilite l'experimentation
et l'iteration rapide.

\subsection{Perspectives d'amelioration}

\begin{itemize}[leftmargin=1.5cm]
  \item \textbf{Modele multilingue} --- Remplacer Wav2Vec2-base par
        \texttt{facebook/wav2vec2-xls-r-300m} (300M parametres, 128 langues)
        pour un support naturel du francais et du darija

  \item \textbf{Fine-tuning partiel} --- Degeler les 2 a 4 dernieres couches
        Transformer de Wav2Vec2 lors de l'entrainement pour une meilleure
        adaptation semantique (necessite GPU)

  \item \textbf{Corpus elargi} --- Passer de 570 a 1\,500 articles Wikipedia
        ($\approx$ 18\,000 chunks) pour couvrir davantage de domaines

  \item \textbf{Hard negative mining} --- Construire des negatifs difficiles
        (documents semantiquement proches mais non pertinents) pour ameliorer
        la discrimination du modele

  \item \textbf{Encoder texte plus puissant} --- Remplacer MiniLM (384D) par
        \texttt{all-mpnet-base-v2} (768D) pour des representations textuelles plus riches

  \item \textbf{Augmentation audio} --- Ajout de bruit de fond, pitch shifting,
        et changement de debit pour la robustesse

  \item \textbf{Deploiement Snowflake} --- Exploiter le pipeline SQL genere
        pour une integration en production dans un data warehouse d'entreprise
\end{itemize}

% ═════════════════════════════════════════════════════════════
\section{Annexe : Structure du projet}
% ═════════════════════════════════════════════════════════════

\begin{lstlisting}[language=bash, caption={Arborescence du projet}]
inpt_audioprocessing/
|-- src/
|   `-- speech_encoder.py        # SpeechEncoder (Wav2Vec2), _load_wav()
|-- training/
|   |-- models.py                # DualEncoderModel (audio + text projections)
|   |-- losses.py                # contrastive_loss (InfoNCE), triplet_loss
|   |-- metrics.py               # Recall@5, Recall@10, MRR
|   |-- data.py                  # Dataset, TextEmbeddingStore, fast collate
|   |-- train_dual_encoder.py    # Script d'entrainement principal
|   |-- precompute_text.py       # Pre-calcul embeddings MiniLM
|   |-- export_text_embeddings.py # Export embeddings 256D post-training
|   `-- validate_handoff.py      # Validation des donnees P1/P2
|-- api/
|   `-- main.py                  # API FastAPI (5 endpoints)
|-- demo/
|   `-- app.py                   # Interface Gradio (2 onglets)
|-- notebooks/
|   |-- P1_Data_Engineer_v3.ipynb
|   |-- P2_Speech_AI_Engineer_v2_FIXED.ipynb
|   `-- S2R_Colab_Kaggle.ipynb   # Notebook cloud training
|-- embeddings/                  # Matrices numpy + manifests CSV
|-- models/                      # Checkpoints entraines
|-- data/                        # Audio + texte + paires
|-- inference.py                 # CLI end-to-end
|-- run_pipeline.bat             # Pipeline Windows
|-- run_pipeline.sh              # Pipeline Linux/macOS
|-- snowflake_pipeline.sql       # Pipeline Snowflake complet
`-- requirements.txt
\end{lstlisting}

% ═════════════════════════════════════════════════════════════
\section{References}
% ═════════════════════════════════════════════════════════════

\begin{enumerate}[leftmargin=1.5cm, label={[\arabic*]}]
  \item Baevski et al. (2020). \textit{wav2vec 2.0: A Framework for Self-Supervised
        Learning of Speech Representations}. NeurIPS 2020.

  \item Conneau et al. (2021). \textit{Unsupervised Cross-lingual Representation
        Learning for Speech Recognition}. Interspeech 2021.

  \item Reimers \& Gurevych (2019). \textit{Sentence-BERT: Sentence Embeddings
        using Siamese BERT-Networks}. EMNLP 2019.

  \item Johnson et al. (2019). \textit{Billion-scale similarity search with GPUs}.
        IEEE Transactions on Big Data.

  \item Panayotov et al. (2015). \textit{LibriSpeech: an ASR corpus based on
        public domain audio books}. ICASSP 2015.

  \item Chen et al. (2020). \textit{A Simple Framework for Contrastive Learning
        of Visual Representations (SimCLR)}. ICML 2020.

  \item Radford et al. (2021). \textit{Learning Transferable Visual Models
        From Natural Language Supervision (CLIP)}. ICML 2021.

  \item HuggingFace. \textit{openslr/librispeech\_asr}.
        \url{https://huggingface.co/datasets/openslr/librispeech_asr}

  \item HuggingFace. \textit{facebook/wav2vec2-base-960h}.
        \url{https://huggingface.co/facebook/wav2vec2-base-960h}

  \item HuggingFace. \textit{sentence-transformers/all-MiniLM-L6-v2}.
        \url{https://huggingface.co/sentence-transformers/all-MiniLM-L6-v2}

  \item FAISS. \textit{Faiss: A library for efficient similarity search}.
        \url{https://github.com/facebookresearch/faiss}

  \item Snowflake Inc. \textit{Snowflake VECTOR type and similarity functions}.
        Snowflake Documentation, 2024.
\end{enumerate}

% ═════════════════════════════════════════════════════════════
\end{document}
% ═════════════════════════════════════════════════════════════
